% Сборка из tex/result: pdflatex -interaction=nonstopmode -halt-on-error full_moments_ridge.tex (дважды).
\documentclass[11pt,a4paper]{article}
\usepackage[T2A]{fontenc}
\usepackage[utf8]{inputenc}
\usepackage[english,russian]{babel}
\usepackage[margin=21mm]{geometry}
\usepackage{amsmath,amssymb,microtype}
\usepackage[hidelinks]{hyperref}
\usepackage{xurl}
\setlength{\parindent}{0pt}
\setlength{\parskip}{4pt}
\setlength{\emergencystretch}{2em}
\newcommand{\R}{\mathbb R}
\newcommand{\norm}[1]{\lVert#1\rVert}
\newcommand{\code}[1]{\nolinkurl{#1}}
\begin{document}
\begin{center}
{\Large\bfseries Полные локальные моменты + ridge}\\[4pt]
{\small Экспериментальный вариант \code{full}; реализация на 25 октября 2026 года}
\end{center}

\paragraph{1. Какие статистики используются.}
Искомый базис $B\in\R^{d\times m}$ удовлетворяет $B^\top B=I_m$.
Внутри одного вызова солвера окна ADP фиксированы. Для весов $w_{ji}\ge0$
определяются
\[
 N_j=\sum_iw_{ji}>0,\qquad a_{ji}=\frac{w_{ji}}{N_j},\qquad
 \pi_j=\frac{N_j}{\sum_kN_k},\qquad
 \mu_j=\sum_i a_{ji}x_i,\quad \bar y_j=\sum_i a_{ji}y_i.
\]
Положим $z_{ji}=x_i-\mu_j$, $\widetilde y_{ji}=y_i-\bar y_j$. Сохраняются
первые локальные моменты по всем $d$ координатам:
\begin{equation}\label{eq:moments}
 C_j=\sum_i a_{ji}z_{ji}z_{ji}^\top\in\R^{d\times d},\qquad
 c_j=\sum_i a_{ji}z_{ji}\widetilde y_{ji}\in\R^d.
\end{equation}
Исходный проекционный критерий использует только
$\Phi_jc_j$ и $\Phi_jC_jB$. Вариант \code{full} строит из \eqref{eq:moments}
полную отбелённую задачу, устраняя потерю компонент при умножении на $\Phi_j$.

\paragraph{2. Отбеливание и вырожденные окна.}
Пусть $C_j=\sum_r\rho_{jr}v_{jr}v_{jr}^\top$. В коде сохраняются моды
$S_j=\{r:\rho_{jr}>10^{-12}\rho_{j,\max}\}$. Обозначим
\[
 \widehat C_j=\sum_{r\in S_j}\rho_{jr}v_{jr}v_{jr}^\top,\qquad
 \Pi_j=\sum_{r\in S_j}v_{jr}v_{jr}^\top,\qquad \widehat c_j=\Pi_jc_j.
\]
Массивы, передаваемые профилировщику, равны
\begin{equation}\label{eq:white}
 U_j=\widehat C_j^{1/2},\qquad
 I_j=\widehat C_j^{\dagger/2}c_j
     =\sum_{r\in S_j}\frac{v_{jr}^\top c_j}{\sqrt{\rho_{jr}}}v_{jr}.
\end{equation}
Поэтому $\norm{I_j-U_jBg}^2$
$=(\widehat c_j-\widehat C_jBg)^\top\widehat C_j^\dagger
(\widehat c_j-\widehat C_jBg)$.
Псевдообратная матрица явно не строится: применяются спектральные множители
из \eqref{eq:white}. Учитываются все сохранённые моды ковариации;
численно отброшенные моды в критерий не входят.

\paragraph{3. Куда добавляется ridge.}
Оптимизируется профиль локальных линейных коэффициентов
\begin{equation}\label{eq:objective}
 \boxed{\displaystyle
 F_\tau(B)=\sum_j\pi_j\min_{g_j\in\R^m}
 \left\{\norm{I_j-U_jBg_j}^2+\tau\norm{g_j}^2\right\},
 \qquad B^\top B=I_m.}
\end{equation}
В исследованной конфигурации для входного базиса $B_{\rm in}$ выбирается
\begin{equation}\label{eq:ridge}
 s=\frac1m\operatorname{median}_j\norm{U_jB_{\rm in}}_F^2,
 \qquad \tau=0.01s.
\end{equation}
Значение $\tau$ фиксируется на весь внутренний вызов. Это статистический
штраф на величину $g_j$. При $s>0$ для каждого пробного $B$ коэффициенты находятся
из независимых систем размера $m\times m$:
\begin{equation}\label{eq:profile}
 \boxed{\displaystyle
 g_j(B)=(B^\top\widehat C_jB+\tau I_m)^{-1}B^\top\widehat c_j.}
\end{equation}
При $\tau>0$ матрица положительно определена даже для вырожденной
$\widehat C_j$. В коде используется решение системы, без явного обращения.
Внешний вес $\pi_j>0$ сокращается в локальном нормальном уравнении.

\paragraph{4. Почему это локальная регрессия.}
Без спектрального усечения $c_j\in\operatorname{range}(C_j)$, и раскрытие
квадратов даёт для любых $B,g_j$ точное равенство
\begin{align}\label{eq:wls}
 &\sum_i a_{ji}(\widetilde y_{ji}-g_j^\top B^\top z_{ji})^2
       +\tau\norm{g_j}^2\nonumber\\
 &\quad=s_{y,j}^2-c_j^\top C_j^\dagger c_j
       +\norm{C_j^{\dagger/2}c_j-C_j^{1/2}Bg_j}^2
       +\tau\norm{g_j}^2,
 \qquad s_{y,j}^2=\sum_i a_{ji}\widetilde y_{ji}^2.
\end{align}
Первые два члена справа не зависят от $B,g_j$. Таким образом, критерий
эквивалентен взвешенной локальной линейной ridge-регрессии с устранёнными
интерсептами. При усечении такое же равенство относится к предикторам
$\Pi_jz_{ji}$. Скаляр $s_{y,j}^2$ для оптимизации \eqref{eq:objective}
не требуется.

\paragraph{5. Как обновляется подпространство.}
После подстановки \eqref{eq:profile} градиент вычисляется по теореме
об огибающей, без дифференцирования решения локальной системы:
\[
 G=2\sum_j\pi_j(\widehat C_jBg_j-\widehat c_j)g_j^\top,
 \qquad \operatorname{grad}F_\tau(B)=(I_d-BB^\top)G.
\]
Критерий инвариантен к поворотам $B\mapsto BQ$, $g_j\mapsto Q^\top g_j$
при $Q^\top Q=I_m$. В проведённых опытах оптимизация реализована через
QR-карту
\[
 B(K)=\operatorname{qf}(B_{\rm in}+B_\perp K),\qquad
 K\in\R^{(d-m)\times m},
\]
и L-BFGS-B с бюджетом до 20 итераций; $B_\perp$ --- ортонормированное
дополнение к $B_{\rm in}$. Между внешними шагами ADP окна и статистики
пересчитываются. Штатный \code{select_step="best"} выбирает внешний шаг
по ошибке локальных средних в центрах.

\paragraph{6. Стоимость и границы.}
Хранение плотных $C_j$ требует $O(Jd^2)$ памяти. После подготовки моментов
одно вычисление профиля и градиента стоит
$O\bigl(J(d^2m+dm^2+m^3)\bigr)$; спектральная подготовка стоит $O(Jd^3)$
на внутренний вызов. Используются локальные линейные модели всех $d$
координат; кривизна функции может давать смещение. Улучшение критерия
и стабилизация систем ridge не гарантируют уменьшения истинной EDR.
Сравнение с обычным Grassmann включает также замену внутреннего оптимизатора.

\smallskip
{\small\textbf{Соответствие коду.}
\code{experiments/coupled_penalties.py}: функции
\code{full_moments}, \code{build}, \code{profile}, \code{solve}, arm \code{full}.
Сбор $C_j,c_j$ и подключение к ADP:
\code{experiments/full_moments_meanmave.py} и
\code{experiments/curvature_penalty_study.py}.}

\paragraph{7. Ускорение без изменения статистического критерия.}
Экспериментальная реализация \code{experiments/full_ridge_batched.py}
профилирует коэффициенты непосредственно по полным моментам. Для каждого
окна достаточно хранить $H_j=\widehat C_j$, $h_j=\widehat c_j$ и
$e_j=c_j^\top\widehat C_j^\dagger c_j$. Обозначим
\[
 T_j=H_jB,\quad v_j=B^\top h_j,\quad
 A_j=B^\top T_j+\tau I_m,\quad g_j=A_j^{-1}v_j.
\]
После раскрытия квадрата и подстановки $A_jg_j=v_j$ получаем
\[
 F_\tau(B)=\sum_j\pi_j(e_j-v_j^\top g_j),\qquad
 G=2\sum_j\pi_j(T_jg_j-h_j)g_j^\top.
\]
Это тот же критерий, включая аддитивную константу. На каждой оценке
профиля требуется только одно умножение $H_jB$; построение и повторное
применение матриц отбеливания исключены. Масштаб \eqref{eq:ridge},
нормировка критерия $\max(1,\sum_j\pi_je_j)$, QR-карта и настройки
L-BFGS-B сохранены. Проекционные моменты по $\Phi_j$ не строятся.
В основной ветке не строится граф и не сохраняются матрицы отбеливания.

Спектральная подготовка выполняется пакетно, как в исходном варианте:
маска, корни и цель сначала вычисляются с теми же формулами; корни нужны
для входного масштаба ridge и затем освобождаются. В полноранговых окнах
используются исходные $C_j,c_j$, в остальных --- их спектральные проекции.
Если входной $\tau\le10^{-8}\max(\max_j\rho_{j,\max},\mathrm{tiny})$,
весь внутренний вызов использует исходный отбелённый профиль. Это решение
фиксируется до line search. Защита нужна вблизи численного ядра ковариации:
при крайне малом ridge алгебраически равные свёртки могут существенно
расходиться в машинной арифметике. В этой резервной ветке сохраняются
матрицы отбеливания и вызывается исходная графовая обвязка.

\paragraph{8. Проверенная, но отклонённая подготовка через Холецкого.}
Первый кандидат \code{experiments/full_ridge_fast.py}
для хорошо обусловленных окон проверял полный ранг через Холецкого.
Пусть $u_j=\max(\norm{C_j}_\infty,\mathrm{tiny})$,
$\delta_j=10^{-7}u_j$. Если для вычисленного фактора $L_j$ матрицы
$C_j-\delta_jI_d$ выполнено
\[
 r_j=\norm{C_j-\delta_jI_d-L_jL_j^\top}_\infty<\delta_j/4,
\]
то в точной арифметике для этих матриц
$\lambda_{\min}(C_j)\ge\delta_j-r_j>10^{-12}u_j$:
все моды проходят исходный порог. Консервативный запас $10^{-7}$
также отправляет плохо обусловленные окна в спектральную ветку.
Для $e_j$ решается система с \emph{исходной} $C_j$, без сдвига.
При неуспехе проверки применяется исходная спектральная маска
\eqref{eq:white}. В этой ветке входной масштаб ridge вычисляется через
сохранённые спектральные корни: след $B^\top H_jB$ может терять точность
вблизи ядра $C_j$. Это проверка в арифметике с плавающей точкой,
не интервальный сертификат. На 36 новых парных полных fit первый кандидат
сохранил EDR с максимальным различием $7.3\cdot10^{-13}$, но медианное
сокращение времени составило лишь $6.7\%$: порог $10\%$ не пройден.
Цикл отдельных факторизаций увеличил время солвера. Его проверочные seed
не использовались; пакетный кандидат исследуется на новых данных.

\paragraph{9. Память и границы ускорения.}
Весовые блоки обрабатываются за один проход. Дополнительное хранение
всех весов $O(Jn)$ заменено буфером $O(bn)$ для размера блока $b$.
Полные моменты по-прежнему занимают $O(Jd^2)$; при спектральном переходе
может потребоваться копия усечённых моментов. Внешний цикл ADP сохраняет
матрицу расстояний $J\times n$, поэтому память \emph{всего} fit остаётся
$O(Jn+Jd^2)$, а не $O(bn+Jd^2)$. Спектральная подготовка имеет прежнюю
кубическую верхнюю оценку. Ускорение является сокращением операций и
временных массивов, а не изменением этих верхних порядков сложности.

\paragraph{10. Проверка на новых парных данных.}
Для пакетного варианта выполнено 144 полных fit: 36 пар на отборе и
36 пар на ранее не использованных проверочных seed. В каждой серии
использованы три функции (additive, quadratic, sine), $m=2,3,5$,
$d=25$, $n=5000$, шум $0.3$ и четыре seed на конфигурацию.
Зафиксированные пороги и независимый аудит пройдены в обеих сериях.
Основные замеры включают одинаковый \code{IndexProfiler} с
\code{tracemalloc}; медианы парных отношений приведены ниже.
\begin{center}
\begin{tabular}{lrr}
 & Отбор & Проверка\\\hline
Время fit, новый/исходный & $0.7245$ & $0.7257$\\
Время процесса со стартом & $0.7801$ & $0.7833$\\
Пиковый RSS процесса & $0.9848$ & $0.9860$\\
Максимальное $|\Delta\mathrm{EDR}|$ & $5.7\cdot10^{-11}$ & $2.6\cdot10^{-12}$
\end{tabular}
\end{center}
В проверке совпали выбранные внешние шаги, последовательности числа
итераций и признаки стационарности. Восстановление при
$\mathrm{EDR}\le0.1$ достигнуто в $16/36$ случаях обоими вариантами.
Резервная отбелённая ветка не потребовалась в $180$ внутренних вызовах
проверочной серии. Это подтверждает ускорение на данной сетке, но не
улучшение статистического восстановления и не всеобщую машинную
эквивалентность для любых вырожденных данных.

Дополнительно на $12$ additive-парах проверочных данных выполнены первые
и повторные fit без сбора аллокаций. Медианы парных отношений времени
равны $0.8088$ и $0.7938$ соответственно: сокращение на $19.1\%$ и
$20.6\%$. Пик RSS процесса в этой дополнительной серии имеет отношение
$0.9554$. Эти повторы не входят в статистические gates и не расширяют
число независимых seed. Отдельный стресс только этапа статистик при
$n=20000,J=600,d=25$ сократил RSS по \code{getrusage}
с $179.5$ до $103.1$ МиБ; полные моменты совпали побитово.
Этот выигрыш памяти одного этапа не является выигрышем памяти всего fit.

\smallskip
{\small Воспроизводимый протокол и результаты: каталоги
\code{verification/full-ridge-fast-20261025} и
\code{verification/full-ridge-batched-20261025}; команда
\texttt{\detokenize{python -m experiments.full_ridge_batched_study run selection}}.
Оптимизация остаётся отдельным экспериментом; штатные настройки ADP
не изменены.}
\end{document}
